\subsection{Event selection}\label{Sec:Mult}

Data used in this analysis were collected with a minimum bias trigger, which requires one or more hits in both V0 scintillator arrays in coincidence with proton beams from both directions. The contamination from beam-induced background is removed offline by using the timing information in the V0 detectors and taking into account the correlation between tracklets and clusters in the SPD detector, as discussed in detail in~\cite{ALICE2}. The primary vertex is reconstructed by correlating hits in the two SPD layers and only events with a primary vertex within $\pm$10\,cm of the nominal interaction point along the beam direction are accepted for this analysis. Due to the fast readout time of the SPD, any contamination from out-of-bunch pile-up is rejected. The contamination from in-bunch pile-up events is removed offline by excluding events with multiple vertices reconstructed in the SPD. Any remaining pile-up will be from collisions which produce little or no particles. Due to the required high-multiplicity event-selection, this will have a negligible impact on the final results presented in this article,

The measurements reported here were performed on minimum-bias triggered events that additionally have at least one charged particle measured in the pseudorapidity interval $|\eta| < 1$ (\INELgtO)~\cite{INEL}, corresponding to about 75\,\% of the total inelastic cross section. Two different multiplicity estimators are used, the total charge deposited in the full coverage of both V0 detectors (V0M), and the number of SPD tracklets within $\etarange 0.8$ (\tracklet). For each of these estimators, the multiplicity is classified as a percentile, where 0\% corresponds to the highest and 100\% to the lowest multiplicity. The high-multiplicity events used throughout this article include the top-1\% (0--1\%) and the top-10\% (0--10\%) multiplicity percentiles, with a minimum of 10 reconstructed charged primary tracks at midrapidity (\etarange 0.8). 

